\chapter{实验与评估}

本章对应当前稿中已经收紧后的 claim 集合。具体来说，\S\ref{sec:layera} 对应主循环稳定性与部分 controlled path，\S\ref{sec:throughput} 对应 AArch64 对齐子集上的吞吐对比，\S\ref{sec:layerb} 与 \S\ref{sec:abcompare} 保留给后续实机补证，\S\ref{sec:layerc} 对应 MIPS 生成器案例。这样的组织方式直接继承了 Layer A/B/C 证据结构，避免把当前还没有数据支撑的主张写得过早。

\placeholderfigure{fig:layers}{[图 6-1 待补]\\Layer A / Layer B / Layer C 证据映射图}{本文评估的三层结构。Layer A 用于开发、sanity 与 matched-subset 对照；Layer B 用于真实硬件补证；Layer C 用于量化新 ISA 接入成本。}

\section{Layer A：稳定性与受控样例}
\label{sec:layera}

本节的核心问题不是“当前是否已经发现新的硬件隐藏指令”，而是“现有框架在受控 Layer A 环境中是否稳定、可审计，并能否通过小规模样例验证分类链路的关键路径”。在这一口径下，当前实验已经足以支撑主循环稳定性和部分 controlled path 的论述。

\begin{table}[t]
  \centering
  \caption{Layer A 稳定性与受控样例实验汇总}
  \label{tab:layera}
  \begin{tabular}{p{1.3cm}p{2.8cm}p{3.0cm}p{5.0cm}}
    \toprule
    Run ID & 平台/后端 & 子集/样例 & 主要结果 \\
    \midrule
    R001 & RISC-V \memcage & \code{0x00000000--0x00000fff} & Tested=4064, Raw H=31, H=0, D=97, T=33, X=32 \\
    R002 & RISC-V \memcage & \code{0x00010000--0x00010fff} & Tested=992，worker crash count 为 0 \\
    R003 & AArch64 \ptracebackend & \code{0x00000000--0x0000ffff} & Tested=65536，D(strict)=65536 \\
    R004 & AArch64 \ptracebackend & legal case & 无异常 artifact \\
    R005 & AArch64 \ptracebackend & illegal case & 运行结束但无 artifact，暴露 taxonomy/logging gap \\
    R006 & AArch64 \ptracebackend & privileged case & D(strict)=1 \\
    R007 & AArch64 \ptracebackend & trap candidates & 运行结束但无 artifact，暴露 exec-fault/trap logging gap \\
    \bottomrule
  \end{tabular}
\end{table}

R001 至 R003 表明，当前框架在声明的 Linux 容器环境下已经能够稳定完成短程扫描并产出完整日志。R004 和 R006 进一步说明，合法样例与 privileged-in-user 样例都能走通当前链路。与此同时，R005 与 R007 也提供了一个同样重要的负结果：illegal-at-test 与 known-disassembly exec-fault 这两类路径虽然已经在受控实验中暴露出来，但当前分类器尚未把它们稳定输出为 artifact。

因此，本文在当前稿中采用更收缩的表达：系统已经能够稳定表达 legal 与 privileged-in-user 两类受控路径，并通过 illegal/trap 样例暴露 taxonomy/logging gap。这个结论比“分类器已完整覆盖 H/D/P/T/X”更弱，但与现有证据是一致的。

\section{Layer A：后端吞吐对比}
\label{sec:throughput}

本节只测试一个收缩后的 claim：在当前 AArch64 对齐子集与当前 Layer A 环境中，\memcage\ 的 median 吞吐高于 \ptracebackend{}。我们不把这个结论外推为“\memcage\ 一般优于 \ptracebackend{}”，因为当前数据只覆盖一个架构、一个对齐窗口和三次重复实验。

\begin{table}[t]
  \centering
  \caption{AArch64 matched subset 上的后端吞吐对比}
  \label{tab:throughput}
  \begin{tabular}{p{2.2cm}p{1.3cm}p{3.3cm}p{1.7cm}p{2.6cm}p{3.0cm}}
    \toprule
    后端 & 重复次数 & 扫描范围 & Tested & median 吞吐 & 备注 \\
    \midrule
    \memcage & 3 & \code{0x00000000--0x0003ffff} & 262144 & 9362.29 insn/s & 三次约为 8456.26、9362.29、9362.29 \\
    \ptracebackend & 3 & 同上 & 262144 & 8192.00 insn/s & 三次约为 7489.83、8192.00、8456.26 \\
    \bottomrule
  \end{tabular}
\end{table}

这组实验有两个值得保留的点。第一，两组实验在 coverage 上保持一致，每次都测试 262144 条编码，并记录到相同数量的 D(strict)；这使得对比建立在相同工作量上，而不是建立在不同异常分布之上。第二，在这一严格对齐前提下，\memcage\ 的 median 吞吐约高 14.3\%。因此，本文采用的最稳妥表述是：在当前 matched subset 上，\memcage\ 表现出稳定但有限的吞吐优势。

\section{Layer B：真实硬件复现}
\label{sec:layerb}

本节保留给后续 RISC-V 板卡实验补证。终稿至少应补入以下信息：板卡型号、SoC 或 stepping、Linux 内核版本、\capstone\ 版本、实验命令、三次复跑结果以及日志归档目录。该部分完成后，本文才能就真实硬件可复现性做出正面陈述。

\placeholdertable{tab:layerb}{Layer B 实机环境与复现实验占位}{[表 6-4 待补]\\建议列：板卡、stepping/SoC、内核、命令、三次复跑结果、日志目录。}

\section{Layer A / Layer B 对照}
\label{sec:abcompare}

本节同样保留为占位。正式补写时，应预先固定对齐子集，例如相同的 \code{-b/-e} 范围或相同 seed，并把差异显式分桶为模拟器差异、解码器差异、特权或模式差异、以及状态依赖差异。核心目标不是追求“完全一致”，而是使观察到的差异可以解释。

\placeholdertable{tab:abcompare}{Layer A / Layer B 差异分桶占位}{[表 6-5 待补]\\建议列：子集定义、Layer A 结果、Layer B 结果、差异类型、解释。}

\section{Layer C：生成器案例研究}
\label{sec:layerc}

Layer C 的问题不是“生成器能否完全替代人工实现”，而是“它是否已经把新 ISA bring-up 从零开始编码，转化为规格输入驱动的可运行基线生成问题”。围绕这个问题，当前 MIPS64 quick smoke 已经给出足够直接的证据。

\begin{table}[t]
  \centering
  \caption{MIPS64 生成器案例的最小量化信息}
  \label{tab:mipscase}
  \begin{tabular}{p{3.5cm}p{2.0cm}p{6.5cm}}
    \toprule
    项目 & 当前值 & 含义 \\
    \midrule
    配置文件行数 & 36 & MIPS64 规格输入规模 \\
    生成后端行数 & 497 & 自动生成的 \code{src/arch\_mips.c} \\
    人工检查项 & 3 & \code{MEMCAGE\_TODO} 中尚未关闭的条目 \\
    quick smoke 时间 & 16.27s & 从调用脚本到得到可运行结果的墙钟时间 \\
    \bottomrule
  \end{tabular}
\end{table}

这组结果说明，生成器已经能够显著降低新 ISA 接入的初始成本，但它还没有把全部工程工作自动化。更准确的说法是：当前生成链路把新架构接入问题压缩成一个小规格输入、一个可运行基线后端，以及一个显式人工检查清单。这种输出形态本身就是一种可复用的工程结论。
